\chapter{Preparation}

\section{Understanding the matrix-based approach}

\textit{The following applications of tropical semirings to dynamic dependency graphs, particularly the sections on 
alternative graph structures (\ref{chap:prep:alternative-graph-structure})
and register-only longest chains (\ref{chap:intro:section:register-only-chains})
are closely adapted from the methodologies established by Chadwick \cite{Chadderz}.}

\subsection{The tropical semiring (max-plus algebra)}

Tropical analysis is a branch of pure mathematics that is the study of the tropical semiring.
There are two variants of the tropical semiring, namely the min-plus and max-plus algebras.
We will use the max-plus convention to apply to \textit{longest} paths in directed acyclic graphs.

 
% Define tropical semiring and operations 

\begin{definition}[Semiring] \label{def:semiring}
  A semiring \cite{Wandsnider2020} is a set $R$ along with two operations, addition (+) and multiplication ($\cdot$), which satisfy the 
  following axioms:

  \begin{enumerate}
    \item Addition is associative and commutative.
    \item Multiplication is associative.
    \item There exists an additive identity, which we usually call $\mathbb{0}$. Additionally, $\mathbb{0}$ annihilates 
      $R$, that is, for all $x \in R$, $\mathbb{0} \cdot x = x \cdot \mathbb{0} = \mathbb{0}$.
    \item There exists a multiplicative identity, which we usually call $\mathbb{1}$.
    \item Multiplication left and right distributes over addition.
  \end{enumerate}
\end{definition}

\begin{definition}[Tropical semiring]
  The tropical semiring \cite{Wandsnider2020}, or max-plus algebra, is the object $(\mathbb{R} \cup \{-\infty\}, \oplus, \odot)$ with operations 
defined as:

$$
\begin{aligned}
  a\oplus b &= \max(a,b) \\
  a\odot b &= a+b
\end{aligned}
$$

\end{definition}

The tropical semiring satisfies the axioms of a semiring from Definition \ref{def:semiring}. It
is not a ring as the max operator does not have an inverse.

% Define matrix multiplication in the context of the tropical semiring

\begin{definition}[Matrix operations in the tropical semiring]
For any $n\times n$ matrices $X = (x_{ij})$ and $Y = (y_{ij})$, the entries of $U = X \oplus Y$ and $V = X \odot Y$
are calculated as \cite{krivulin2012maxplusalgebraapproachmodelling}

$$
u_{ij} = x_{ij} \oplus y_{ij}, \text{ and } v_{ij} = \displaystyle\sum_{k=1}^{n} \! {}_{\oplus} x_{ik} \odot y_{kj}
$$

\end{definition}

The matrix operations $\oplus$ and $\odot$ are associative owing to the commutativity and associativity of the underlying 
addition operator, and that the multiplication operator distributes over addition.

\subsection{Longest paths in Directed Acyclic Graphs}


% Show how this applies to the longest-path problem

As established previously, finding the length of the critical path in a program amounts to finding the longest path
in its DDG. It is noted that the max-plus version of the tropical semiring is applicable to finding the longest path
in a graph.

One algorithm to find the longest path between nodes $A$ and $B$ in a graph is to recursively take the 
maximum of all path lengths from all neighbours of $A$ to $B$, plus the cost of the edge from $A$ to the chosen
neighbour. We can see that this amounts to a Bellman recursion equation \cite{Bellman:DynamicProgramming}
involving only the max and plus operators: 

$$
v(A,B) = 
\begin{cases}
0 &\text{ if } A = B \\
\max_{\forall C . A \to C} \{\mathrm{weight}(A \to C) + v(C,B))\}, &\text{ otherwise } \\
\end{cases}
$$

We can observe that, if we have an adjacency matrix (where the cells in that matrix represent the cost of the edge
between the pair of nodes in that row/column, and $-\infty$ if there is no edge), performing tropical semiring 
matrix multiplication of this matrix with itself corresponds to one step of our longest-path algorithm.

Consider an adjacency matrix $X^{2} = X \odot X$, and denote its entries by $x_{ij}^{(2)}$. We have that $x_{ij}^{(2)} = p$
if and only if there exists at least one path of length $p$ from node $i$ to node $j$ in the corresponding graph,
consisting of at most two edges, and that path is the longest such path. This extends to any integer exponent $q$.

\begin{theorem}[Matrix exponentiation results in longest paths] \label{thm:mat-exp-longest-paths}

For any integer $q>0$, the matrix
$X^{q}$ has the entry $x_{ij}^{(q)} = p$ only if there is a path of length $p$ with at most $q$ edges from $i$ to $j$,
and that path is the longest such path. \cite{krivulin2012maxplusalgebraapproachmodelling}.

\begin{proof}
  We proceed by induction on $q$, the maximum number of edges in the path.

  \textbf{Base case ($q = 1$):} Consider the matrix $X^1=X$. By the definition of an adjacency matrix, the entry  
  $x_{ij}^{(1)}$ is exactly the weight of the direct edge from node $i$ to node $j$. If no such edge exists, then 
  $x_{ij}^{(1)} = -\infty$. Thus, the base case holds: the entries represent their longest paths consisting of at most 
  one edge.

  \textbf{Inductive step:} Assume the claim holds for $q$. That is the matrix $X^{q}$ has entries $x_{ij}^{q}$ 
  representing the longest path from $i$ to $j$ using at most $q$ edges.
  We show that the claim holds for $q+1$. By definition, $X^{q+1} = X^{q} \odot X$
  The entry $x_{ij}^{(q+1)}$ is calculated as:
  $$x_{ij}^{(q+1)} = \sum_{k=1}^{n} \! {}_{\oplus} (x_{ik}^{(q)}\odot x_{kj})$$
  Translating this to the original operators:
  $$
  x_{ij}^{(q+1)} = \max_{1 \leq k \leq n} (x_{ik}^{(q)} + x_{kj})
  $$
  This lines up perfectly with the dynamic programming approach to finding the longest path. To find a longest path from
  node $i$ to node $j$, it checks every intermediate node $k$. It takes the longest known path from $i$ to $k$ of at most 
  $q$ edges (by the inductive hypothesis), and adds the direct edge cost from $k$ to $j$. By taking the maximum over all 
  intermediate nodes $k$, we find the absolute longest path from $i$ to $j$ of at most $q+1$ edges. Thus the claim holds 
  for all integers $q>0$.
\end{proof}

So, an adjacency matrix raised to the power infinity will converge to the longest distance matrix for that whole graph.
  
\end{theorem}


For example, consider the graph:
\begin{center}
\begin{tikzpicture}[
    node distance=1cm and 2cm, % vertical and horizontal spacing
    every node/.style={font=\Large} % makes letters look like your image
]
    % Place the nodes
    \node (a) {a};
    \node (b) [right=of a] {b};
    \node (c) [below=of b] {c};
    \node (d) [right=of c] {d};

    % Draw the edges (undirected lines) with labels
    \draw [->] (a) -- node[above] {1} (b);
    \draw [->] (a) -- node[below left] {3} (c);
    \draw [->] (b) -- node[right] {1} (c);
    \draw [->] (c) -- node[above] {1} (d);
\end{tikzpicture}
\end{center}

The corresponding adjacency matrix would be:
$$
X = 
\begin{pmatrix}
  0 & 1 & 3 & -\infty \\ 
  -\infty & 0 & 1 & -\infty \\ 
  -\infty & -\infty & 0 & 1 \\
  -\infty & -\infty & -\infty & 0 \\
\end{pmatrix}
$$

So, in the tropical semiring:

$$
\begin{aligned}
  X^{2} &= 
  \begin{pmatrix}
    0
  \end{pmatrix} \\ 
  X^{3} &= 
  \begin{pmatrix}
    0
  \end{pmatrix}
\end{aligned}
$$

\subsection{Forming the DDG and register-only longest chains} \label{chap:intro:section:register-only-chains}

TODO

\subsection{Alternative graph structure} \label{chap:prep:alternative-graph-structure}

% Show how we arrange the DDG such that we can have a matrix representing dependencies of an instr.

Dynamic dependency graphs typically exhibit a layered, sequential structure - they are more long than wide. While a 
standard $N \times N$ adjacency matrix (where $N$ is the total number of dynamically executed instructions) could 
map the graph, it is inefficient.
This is because, in the DDG, edges \textit{only} 
exist between adjacent columns. Therefore, the full adjacency matrix $X$ is mostly filled with $-\infty$.
Instead, the DDG can be modelled as a sequence of adjacent columns. 
Consider an arbitrary graph with the property of being more long than wide:

\begin{center}
\begin{tikzpicture}[
    node distance=0.7cm and 2cm, % vertical and horizontal spacing
    every node/.style={font=\Large} % makes letters look like your image
]
    % Place the nodes
    \node (a) {a};
    \node (b) [below=of a] {b};
    \node (c) [right=of a] {c};
    \node (d) [below=of c] {d};
    \node (e) [right=of c] {e};
    \node (g) [right=of e] {g};
    \node (f) [below=of e] {f};
    \node (h) [below=of g] {h};

    % Draw the edges (undirected lines) with labels
    \draw [->] (a) -- node[above] {1} (c);
    \draw [->] (a) -- node[below] {4} (d);
    \draw [->] (c) -- node[above] {1} (e);
    \draw [->] (e) -- node[above] {1} (g);
    \draw [->] (f) -- node[above] {2} (g);
    \draw [->] (f) -- node[below] {2} (h);
\end{tikzpicture}
\end{center}

Whilst we could use the adjacency matrix representation, the structure lends itself to another representation for 
longest path problems. First of all, we modify the graph such that the longest path will always connected to the 
rightmost part of the graph, by adding extra nodes and cost-0 edges, and so also each path is connected to the first 
column similarly.

To utilize column-wise matrices for longest-path calculations, the graph must first be padded. We introduce a set of 
'dummy' nodes to the first and last columns, connected by 0-cost edges. This guarantees that any local path within the 
interior of the graph is extended to touch both the initial and final columns.

\begin{center}
\begin{tikzpicture}[
    node distance=0.7cm and 2cm, % vertical and horizontal spacing
    every node/.style={font=\Large} % makes letters look like your image
]
    % Place the nodes
    \node (a) {a};
    \node (b) [below=of a] {b};
    \node (c) [right=of a] {c};
    \node (d) [below=of c] {d};
    \node (e) [right=of c] {e};
    \node (g) [right=of e] {g};
    \node (f) [below=of e] {f};
    \node (h) [below=of g] {h};

    \node (p1) [above=of a, font=\small] {$p_1$};
    \node (p2) [above=of c, font=\small] {$p_2$};
    \node (p3) [above=of e, font=\small] {$p_3$};
    \node (p4) [above=of g, font=\small] {$p_4$};

    \node (m1) [below=of b, font=\small] {$m_1$};
    \node (m2) [below=of d, font=\small] {$m_2$};
    \node (m3) [below=of f, font=\small] {$m_3$};
    \node (m4) [below=of h, font=\small] {$m_4$};

    % Draw the edges with labels
    \draw [->] (a) -- node[above] {1} (c);
    \draw [->] (a) -- node[below] {4} (d);
    \draw [->] (c) -- node[above] {1} (e);
    \draw [->] (e) -- node[above] {1} (g);
    \draw [->] (f) -- node[above] {2} (g);
    \draw [->] (f) -- node[below] {2} (h);

    % draw the zero-cost edges
    \draw [->, dashed] (p1) -- (p2);
    \draw [->, dashed] (p1) -- (c);
    \draw [->, dashed] (p1) -- (d);

    \draw [->, dashed] (p2) -- (p3);
    \draw [->, dashed] (p2) -- (e);
    \draw [->, dashed] (p2) -- (f);

    \draw [->, dashed] (p3) -- (p4);
    \draw [->, dashed] (p3) -- (g);
    \draw [->, dashed] (p3) -- (h);

    \draw [->, dashed] (a) -- (m2);
    \draw [->, dashed] (b) -- (m2);
    \draw [->, dashed] (m1) -- (m2);

    \draw [->, dashed] (c) -- (m3);
    \draw [->, dashed] (d) -- (m3);
    \draw [->, dashed] (m2) -- (m3);

    \draw [->, dashed] (e) -- (m4);
    \draw [->, dashed] (f) -- (m4);
    \draw [->, dashed] (m3) -- (m4);
\end{tikzpicture}

{\tiny dashed edges are 0 cost}
\end{center}

Now we can represent the graph with one matrix per adjacent-column pair, encoding the connectivity between those two
columns. The connectivity between the first two columns is represented as:

\[
  X_{1,2} = 
  \begin{pmatrix}
  0 & 0 & 0 & -\infty \\
  -\infty & 1 & 4 & 0 \\
  -\infty & -\infty & -\infty & 0 \\
  -\infty & -\infty & -\infty & 0 
  \end{pmatrix}
\]

where the $m$th row and $n$th column of this matrix represent the connection from the $m$th node in the first column of 
the graph to the $n$ node in the second column of the graph.

The remaining two column pairs are represented as:

\[
  \begin{aligned}
    X_{2,3} &= 
    \begin{pmatrix}
    0 & 0 & 0 & -\infty \\
    -\infty & 1 & -\infty & 0 \\
    -\infty & -\infty & -\infty & 0 \\
    -\infty & -\infty & -\infty & 0 
    \end{pmatrix} \\
    X_{3,4} &= 
    \begin{pmatrix}
    0 & 0 & 0 & -\infty \\
    -\infty & 1 & -\infty & 0 \\
    -\infty & 2 & 2 & 0 \\
    -\infty & -\infty & -\infty & 0 
    \end{pmatrix}\\
  \end{aligned}
\]

These matrices are each $4 \times 4$ whereas a full adjacency matrix for the original graph would be $7 \times 7$.
We can see that the column-wise matrices will use less storage than the adjacency matrix (48 integers vs 49 integers),
and this alternative becomes increasingly favourable as the graph gets longer.


In the tropical semiring, computing the matrix product

\[
  X_{3,4} \odot X_{2,3} \odot X_{1,2} = \begin{pmatrix}
  0 & 2 & 2 & 1 \\
  -\infty & 3 & -\infty & 4 \\
  -\infty & -\infty & -\infty & 0 \\
  -\infty & -\infty & -\infty & 0 
\end{pmatrix}
\]
yields a matrix representing the longest path from the $m$th node in the first column the $n$th node in the final 
column of the graph. Because we have ensured every path in the graph connects to the first and last column, the largest 
individual value in this matrix is the length of the longest path in the graph, namely the path of length 4 connecting 
from $a$ in the first column, to $m_4$ in the final column, via $d$ and $m_3$.

The correctness of the column-wise approach for calculating longest paths in graphs is a corollary that follows directly 
from Theorem \ref{thm:mat-exp-longest-paths}.

\begin{corollary}[Correctness of the column-wise formulation]
  In a standard adjacency matrix $X$, computing $X^{q}$ checks every possible intermediate node 
  $k$ in the entire graph to find paths of length $q$, as shown by Theorem \ref{thm:mat-exp-longest-paths}.

  In the DDG, when we multiply adjacent column matrices (e.g., $X_{2,3} \odot X_{1,2}$),
  we are simply performing the exact operations for finding the longest path between columns 1 and 3,
  but restricted only to the intermediate nodes $k$ located in column $2$.
  This is valid as edges in the DDG only exist between adjacent columns.

  Because the padding of 0 cost edges ensures every path spans from the first column to the final column, multiplying 
  all consecutive column matrices sequentially computes the maximum path length across the entire graph.
\end{corollary}

This defines the shape of the data will be working on. 

\section{Datasets}
